% ==============================================================
\section{はじめに}\label{sec:intro}
% ==============================================================

\subsection{動機と応用背景}

大規模小売チェーンのデータアナリストは、
「どの商品の組み合わせが、どの店舗で、いつ集中的に購買されるか」
を知りたいと考える。
例えば、100 店舗 $\times$ 52 週の POS データから
「ビールとポテトチップスは週 46--52（年末商戦）の都市部店舗（店舗 0--9）で
集中購買される」という知見は、
地域特化の在庫計画やプロモーション施策に直結する。
同様に、疫学では特定地域・特定時期に集中する症状の共起パターンが
介入ターゲットとなり、
物流では地域・時間帯別の出荷量急増パターンが
サービス最適化を支援する。

このような\emph{時空間局所的な密集パターン}を発見するには、
トランザクションデータの時刻情報と空間情報
（店舗 ID・地域コード・GPS 座標等）を
同時に扱うマイニング手法が必要である。

\subsection{同時出力の必要性と既存手法の限界}

上記の応用において分析者が真に必要とする出力は、
(i)~\emph{何が}密集しているか（アイテムセット $P$）と
(ii)~\emph{いつ・どこで}密集しているか（$d$ 次元密集領域）の
組 $(P, R(P))$ である。
パターン $P$ だけでは「どの店舗にいつ追加発注すべきか」が不明であり、
領域 $R$ だけでは「何を発注すべきか」が特定できない。
両者が揃って初めて、地域・時期に特化した施策を立案できる。
しかし、既存手法はいずれもこの両方を同時に提供できない。

\paragraph{パターンマイニング系手法.}
頻出パターンマイニング~\cite{agrawal1994apriori,han2000mining}は
グローバルサポートに基づくため局所的密集パターンを見逃す。
Apriori-Window~\cite{aprioriwindow}は時間軸のみの1次元設計であり
空間情報を活用できない。
これらは「何が密集しているか」には答えられるが
「どこで密集しているか」には答えられない。

\paragraph{領域抽出系手法.}
密部分テンソルマイニング（M-Zoom~\cite{shin2018mzoom}等）は
直積ブロック制約により任意形状の密集領域にフィットできず、
パターン列挙も不正確である（\S\ref{sec:related}で詳述）。
空間統計手法~\cite{anselin1995lisa,ord1995getis}は時間次元を扱えない。
スキャン統計~\cite{kulldorff1997spatial}は単一変量限定である。
ST-DBSCAN~\cite{birant2007stdbscan}は点近傍密度に基づくため
背景密度と区別された局所的密集を検出できず、
反単調性も満たさないため Apriori 枝刈りと統合できない。

\subsection{本研究の貢献}

本稿では、Apriori-Window~\cite{aprioriwindow}を
任意の $d$ 次元グリッド
（時間 + $d{-}1$ 個の空間・カテゴリ次元）に拡張した
\textbf{多次元 Apriori-Window}を提案する（\S\ref{sec:method}）。
提案手法の特長を一言で表すと、
\textbf{アナリストが「いつ・どこで・何が売れるか」を偽陽性ゼロで抽出できる}
ことである。

本研究の技術的本質は、
1次元 Apriori-Window~\cite{aprioriwindow}の密集区間計算を、
固定空間開始点ごとの時間多重集合へ適用する形で
$d$ 次元へ拡張した点にある。
すなわち、多次元化しても本質的に解いているのは
\emph{重複を許す1次元 Apriori-Window} である。

\begin{enumerate}
\item \textbf{固定空間開始点ごとの 1D 還元}：
各空間開始点 $\bm{s}_{\mathrm{sp}}$ ごとに、
その空間窓に入るパターンの出現を時刻の多重集合
$L_P(\bm{s}_{\mathrm{sp}})$ へ潰す。
これにより $d$ 次元問題は、
空間開始点をキーとする複数の1次元問題に分解される。

\item \textbf{duplicate-aware $t$ 軸ストライド}：
2次元以上では同一時刻の重複が生じる。
1点ずつ dense start を保持するのではなく、
現在窓だけで dense が保証される left の連続ブロックをまとめて処理し、
密集区間として出力する（Lemma~\ref{lem:stride}）。

\item \textbf{singleton interval map の交差による Apriori 枝刈り}：
各アイテムの密集領域を
空間開始点ごとの時間区間列として保持し、
候補パターンではそれらを交差して候補区間を作る。
これにより $k \geq 2$ でも1次元と同型の区間ベース枝刈りが適用できる
（Lemma~\ref{lem:apriori}）。
\end{enumerate}
